%%This is a standard LaTeX2e article document template. personal version 12/5/200%%
\documentclass[12pt,twoside,fullpage]{amsart}
%%%%%%%%%%%%%%%%%%%%%%%%%%%%%%%packages%%%%%%%%%%%%%%%%%%%%%%%%%%%%%%%%%%%%%%%%%%%%%%%%%%%%%%%%%%
%\usepackage{epsf}
%        \usepackage{graphicx}
%        \usepackage[scanall]{psfrag}
%\pagestyle{empty}

\usepackage{latexsym}
\usepackage{amssymb}
\usepackage{amsfonts}
\usepackage{amstext}
\usepackage{fullpage}
\setlength{\topmargin}{-0.75in}
\setlength{\textheight}{10in}
%\setlength{\textwidth}{7in}
%\setlength{\leftmargin}{-0.5in}
%\setlength{\headsep}{0in}          %%%
%\setlength{\headheight}{0in}
%\newcommand{\tld}{\widetilde }
%\newcommand{\bR}{\mathbb{R}}
\renewcommand{\u}{{\bf{u}}}
\renewcommand{\v}{{\bf{v}}}
\newcommand{\w}{{\bf{w}}}
\renewcommand{\t}{{\bf{t}}}
\newcommand{\x}{{\bf{x}}}
\newcommand{\y}{{\bf{y}}}
\newcommand{\z}{{\bf{z}}}
\newcommand{\0}{{\bf{0}}}
\pagestyle{empty}
\newcommand{\bR}{\mathbb{R}}
\newcommand{\bN}{\mathbb{N}}
\newcommand{\bQ}{\mathbb{Q}}
\newcommand{\bZ}{\mathbb{Z}}
\newcommand{\nsubset}{\not\subset}

\newtheorem{theorem}{Theorem}[section]
\newtheorem{lemma}[theorem]{Lemma}
\newtheorem{proposition}[theorem]{Proposition}
\newtheorem{corollary}[theorem]{Corollary}
\newtheorem{conjecture}[theorem]{Conjecture}
\newtheorem{excercise}[theorem]{Exercise}

\theoremstyle{definition}
\newtheorem{definition}[theorem]{Definition}
\newtheorem{example}[theorem]{Example}
\newtheorem{xca}[theorem]{Exercise}
\newtheorem{exercise}[theorem]{Exercise}

\theoremstyle{remark}
\newtheorem{remark}[theorem]{Remark}
\newtheorem{hint}[theorem]{Hint}
\newtheorem{question}[theorem]{Question}


\begin{document}
\title{Worksheet on linear transformations}
\author{Math 186--1}
\maketitle

\begin{definition}Given a function $T$ (which takes vectors as input, and outputs vectors), we say that $T$ is a \emph{linear transformation} if the following two properties hold.
    \begin{itemize}
    \item[(\#1)]  For any two vectors ${\bf v}$ and ${\bf v'}$, we always have $T({\bf v + v'}) = T({\bf v}) + T({\bf v'})$.  In other words, adding the vectors and then using the function gives the same result as using the function, then adding the two outputs.
    \item[(\#2)]  For any ${\bf v}$ and any real number $c$, we have $T(c {\bf v}) = c T({\bf v})$.  In other words, if you scale the vector and then apply the function, you get the same result as if you'd applied the function and then scaled the output.
    \end{itemize}
\end{definition}
\vskip 12pt
\hskip -12pt
{\bf 1.}  Suppose that $T : \bR^2 \to \bR^2$ is a linear transformation ($\bR^2$ we view as the set of vectors in the plane, and it includes the zero vector).  Fix $\u$ and $\v$ to be a basis.  For all other vectors we write $\x = a_x \u + b_x \v$ because $\u, \v$ are a basis.\footnote{Other common expressions for $\x$ include $(a_\x, b_\x) = \left[ \begin{array}{c} a_\x\\b_\x \end{array}\right]$.}  Show that $T(\x)$, for any $\x$, is completely determined by $T(\u)$ and $T(\v)$.  
\vskip 175pt
\hskip -12pt
{\bf 2.}  Fix a basis $\u$, $\v$ for $\bR^2$.  Show that for every other pair of vectors $\x, \y \in \bR^2$, that there exists a linear transformation $T$ such that $T(\u) = \x$ and $T(\v) = \y$.  
\newpage
\hskip -12pt
{\bf 3.}  Write down matrix representations of the following linear transformations also explain as well as you can what this linear transformation does geometrically.  Fix a basis $\u, \v$ for $\bR^2$ and a basis $\x, \y, \z$ for $\bR^3$.  
\begin{itemize}
 \item[(a)]  The map $T : \bR^2 \to \bR^2$ defined in the following way.  $T(\u) = \v$ and $T(\v) = \u$.
\vskip 140pt
 \item[(b)]  The map $T : \bR^2 \to \bR^2$ defined in the following way.  $T(\u) = \v$ and $T(\v) = -\u$.    
\vskip 140pt
 \item[(c)]  The map $T : \bR^2 \to \bR^2$ defined in the following way.  $T(\u) = \u$ and $T(\v) = \0$.   
\vskip 140pt
 \item[(d)]  The map $T : \bR^2 \to \bR^3$ defined in the following way.  $T(\u) = \x$ and $T(\v) = \y + \z$.  
\vskip 140pt
 \item[(e)]  The map $T : \bR^3 \to \bR^2$ defined in the following way.  $T(\x) = \u$ and $T(\y) = \v$ and $T(\z) = \0$.
\end{itemize}

\end{document}
